\section{Introduction}
There is growing concern regarding the phenomenon of ``AI sycophancy'' among researchers and users alike. A growing number of case studies have illustrated the potential catastrophic harms caused by excessive validation from AI systems [CITE], including hypotheses about potential links between sycophancy and an emerging phenomenon known as ``AI psychosis''. The critical nature of this topic necessitates careful empirical evaluations of these diverse phenomena and their downstream impacts. However, many AI sycophancy works are focused on a small subset of sycophantic behaviors constrained within a particular setting. There is a lack of data on the broader impacts of different forms of AI sycophancy, particularly the medium-to-long-term impacts. Further, there is a lack of consensus as to what ``counts'' as AI sycophancy among researchers, and no standardized definition of this term, which can result in broad generalizations being made about the nature of AI sycophancy and certain behaviors being overlooked in favor of those that are easier to evaluate. \citet{ye2026countsaisycophancytaxonomy} recently introduced a unified taxonomy for categorizing and discussing different types of AI sycophancy, which is a promising path towards standardizing the language we use to discuss AI sycophancy, and identified disagreements among researchers regarding the types of behaviors that should be considered sycophantic, referring to AI sycophancy as a \emph{fragmented construct}.

To address these issues in the field, we propose a unified, systematic, and normative approach to characterizing and evaluating AI sycophancy. We argue for a simple definition of AI sycophancy that aligns with the behaviorist nature of AI sycophancy research and promotes explorations of downstream harms of people-pleasing behaviors in AI systems. This definition can be formalized as \emph{user-approval-driven norm displacement}: a model behavior that accommodates a user’s belief, affect, identity, goal, or framing in order to preserve approval or relational ease, while reducing fidelity to a context-relevant norm such as truth, calibration, autonomy, care, fairness, task quality, or institutional integrity. We believe that this definition, and the accompanying framework we introduce, promote open investigation and documentation of the diverse behaviors in AI systems that optimize short-term user preferences over long-term outcomes. 

Our contributions can be listed as follows:

\begin{enumerate}
    \item Propose an empirical approach to characterizing AI sycophancy, which can resolve disagreements among researchers as to the types of behaviors that constitute sycophancy.
    \item Introduce a step-by-step framework for characterizing and evaluating AI sycophancy that centers impact and promotes the evaluation of a diverse set of behaviors
    \item Provide a roadmap for practitioners, which includes recommended best practices and gaps in the literature to inform future research
\end{enumerate}

% \begin{enumerate}
%     \item Motivation: findings related to people-pleasing in LLMs and case studies showing the harms of sycophancy
%     \item Trends in AI research, news coverage, and online discourse
%     \item Position: how researchers measure AI sycophancy often diverges from users' perceptions of AI sycophancy and major concerns about it.
%     \item Position: we need a more impact-centered framework for characterizing and evaluating AI sycophancy
%     \begin{enumerate}
%         \item Can help us better understand the dangers of this phenomenon, and which sycophantic behaviors causes the most harm, so we can target our mitigation strategies
%         \item Can also help us distinguish sycophancy from things like personalization and empathetic behavior - the line between adjusting for human preferences and excessive people-pleasing could be drawn based on the negative impacts over time
%     \end{enumerate}
%     \item One way to distinguish sycophancy from normal preference optimization, in the case of a lack of ground truth, could be by studying the impacts of this behavior (characterize ``excessive'')
%     \item Roadmap for the rest of the paper: discuss existing literature, argue for the need for more information regarding impacts of AI sycophancy, and enumerate recommendations for future work in this space
%     \end{enumerate}

    